\subsection{指数与阶}

定理 6. —— 设 K 是一个交换域，并设 A 是 K 上的一个有限次数中心单代数。设 L 是 K 的一个有限次数可分扩张，它是代数 A 的一个分裂域。则 \([L:K][A]\) 在 \(\mathrm{Br}(K)\) 中为零。

存在 L 的一个扩张 M，它是 K 的一个有限次数伽罗瓦扩张（V, §10, No.~1, 第 57 页，命题 2）。A 的类 {[}A{]}（在 K 的布饶尔群中的）属于子群 Br(M/K)。设 G 是 M 在 K 上的伽罗瓦群，并设 \(\alpha\) 是 {[}A{]} 在 \(H^{2}(G,M^{*})\) 中的像（VIII, 第 321 页，推论）。设 H 是 M 在 L 上的伽罗瓦群。则 H 是 G 中指标为 {[}L:K{]} 的子群（V, §10, No.~7, 第 68 页，推论 5）。由于 \(\operatorname{Res}_{\mathrm{H}}^{\mathrm{G}}(\alpha)=\Phi_{\mathrm{M/L}}(\mathrm{A}_{(\mathrm{L})})\)（命题 14），我们有 \(\operatorname{Res}_{\mathrm{H}}^{\mathrm{G}}(\alpha)=0\)。由 VIII, 第 303 页的命题 8 推出 {[}L:K{]} \(\alpha=0\)，从而 {[}L:K{]}{[}A{]}=0。

设 K 是一个交换域，并设 A 是 K 上的一个有限次数中心单代数。则 A 同构于形如 \(\mathbf{M}_{n}(\mathrm{D})\) 的代数，其中 D 是一个中心为 K 的体，且 \([A]=[D]\)（在 \(\mathrm{Br}(K)\) 中）。D 的约化次数只依赖于 A。我们把这个约化次数称为 A 的指数。A 的指数整除 A 的约化次数。A 的阶是 A 的类在 K 的布饶尔群中的阶。

推论 1. —— 交换域上的有限次数中心单代数的阶整除其指数。

只需对一个体 D（在其中心 K 上的次数为有限的）证明这一点。设 L 是 D 的一个极大交换子域，它是 K 的一个可分扩张；我们有 \([D:K]=[L:K]^{2}\)（VIII, 第 265 页，推论 2, b) 与 c)）。于是 K 的扩张 L 是代数 D 的一个分裂域（VIII, 第 281 页，命题 5），且 \([L:K]\) 与 D 的约化次数一致。然后我们应用定理 6。

推论 2. —— 设 K 是一个交换域，并设 A 是 K 上的一个有限次数中心单代数。设 p 是一个素数。若 p 整除 A 的指数，则 p 整除 A 的阶。

我们假设素数 \(p\) 不整除 A 的阶，并证明它不整除 A 的指数。只需在 A 是一个体的情况下证明这一结果。设 L 是 K 的一个分裂 A 的有限次数伽罗瓦扩张。设 G 是 L 在 K 上的伽罗瓦群，并设 H 是 G 的一个西罗 \(p\)-子群（I, §6, No.~6, 第 78 页）。我们用 \(\mathrm{L}' = \mathrm{L}^{\mathrm{H}}\) 表示 L 中 H 的不变量所成的子域。\(\mathrm{A}_{(\mathrm{L}')}\) 的阶整除 A 的阶，从而与 \(p\) 互素。由定理 6，我们有 \([\mathrm{L}:\mathrm{L}'][\mathrm{A}_{(\mathrm{L}')}] = 0\)（在 \(\mathrm{L}'\) 的布饶尔群中）。由此推出 \([\mathrm{A}_{(\mathrm{L}')}] = 0\)，且体 \(\mathrm{L}'\) 分裂 A。然后我们应用 VIII, 第 283 页的推论 2；由此推出 A 的指数整除 \([\mathrm{L}':\mathrm{K}]\)，从而不被 \(p\) 整除。

推论 3. —— 设 p 是一个素数，并设 K 是一个特征为 p 的完全域。设 A 是 K 上的一个有限次数中心单代数。则 p 不整除 A 的指数。

我们证明 K 的布饶尔群不含阶为 p 的元素。K 的每个有限次数伽罗瓦扩张 M 都是完全域（V, §7, No.~1, 第 36 页，命题 2）。于是取 p 次幂是群 \(M^{*}\) 的一个自同构。由此推出，用 p 作乘法是群 \(\mathrm{H}^{2}(\mathrm{Gal}(\mathrm{M}/\mathrm{K}),\mathrm{M}^{*})\) 的自同构，该群同构于 \(\mathrm{Br}(\mathrm{M}/\mathrm{K})\)。

因此，{[}A{]} 的阶与 \(p\) 互素，且由推论 2，其指数不被 \(p\) 整除。

注. —— 通过考虑四元数代数的张量积，可以构造以 K 为中心、阶为 2 且指数任意大的体（参看 VIII, 第 371 页，习题 7 与 8）。

*反之，若 K 是 \(p\)-进域或有限域上形式幂级数域的一个有限伽罗瓦扩张，则 K 上的有限次数中心单代数的阶等于其指数（VIII, 第 332 页，习题 17, e)）。*

